
\documentclass{article}
\usepackage{colortbl}
\usepackage{makecell}
\usepackage{multirow}
\usepackage{supertabular}

\begin{document}

\newcounter{utterance}

\centering \large Interaction Transcript for game `cladder', experiment `full\_v1.5\_default', episode 1290 with Qwen3.5{-}4B.
\vspace{24pt}

{ \footnotesize  \setcounter{utterance}{1}
\setlength{\tabcolsep}{0pt}
\begin{supertabular}{c@{$\;$}|p{.15\linewidth}@{}p{.15\linewidth}p{.15\linewidth}p{.15\linewidth}p{.15\linewidth}p{.15\linewidth}}
    \# & \multicolumn{2}{c}{Player} && \multicolumn{2}{c}{Game Master} \\
    \hline

    \theutterance \stepcounter{utterance}  
    & & & \multicolumn{4}{p{0.6\linewidth}}{
        \cellcolor[rgb]{0.9,0.9,0.9}{
            \makecell[{{p{\linewidth}}}]{
                \texttt{\tiny{[P1$\langle$GM]}}
                \texttt{You are an expert in causal inference. The following question is not a typical commonsense query, but rather a meticulously designed question created by a professor specializing in causal inference, intended to assess the students' mastery of the course content.} \\
\texttt{Imagine a self{-}contained, hypothetical world with only the following conditions, and without any unmentioned factors or causal relationships: Gender has a direct effect on smoking and lung cancer. Smoking has a direct effect on tar deposit. Tar deposit has a direct effect on lung cancer. Gender is unobserved. For nonsmokers, the probability of high tar deposit is 36\%. For smokers, the probability of high tar deposit is 84\%. For nonsmokers and with no tar deposit, the probability of lung cancer is 12\%. For nonsmokers and with high tar deposit, the probability of lung cancer is 44\%. For smokers and with no tar deposit, the probability of lung cancer is 15\%. For smokers and with high tar deposit, the probability of lung cancer is 50\%. The overall probability of smoking is 67\%. Will smoking increase the chance of lung cancer?} \\
\texttt{Start your answer with 'Yes' or 'No', followed by additional reasoning or evidence to support your explanation.} \\
            }
        }
    }
     \\ \\

    \theutterance \stepcounter{utterance}  
    & \multicolumn{4}{p{0.6\linewidth}}{
        \cellcolor[rgb]{0.561,0.737,0.561}{
            \makecell[{{p{\linewidth}}}]{
                \texttt{\tiny{[P1$\rangle$GM]}}
                \texttt{Yes} \\
\\ 
\texttt{**Reasoning:**} \\
\\ 
\texttt{To determine whether smoking increases the chance of lung cancer, we must compare the counterfactual probability of lung cancer for a smoker (\$P(LC | \textbackslash{}text\{Smoker\})\$) with the probability of lung cancer for a nonsmoker (\$P(LC | \textbackslash{}text\{Nonsmoker\})\$). In causal inference, this comparison isolates the effect of the treatment (smoking) by holding the distribution of confounders (in this specific hypothetical world, Gender) constant or by comparing the potential outcomes directly.} \\
\\ 
\texttt{Since the problem states that "Gender is unobserved" but also explicitly defines the causal pathways and probabilities conditional on smoking status without mentioning any other unobserved confounders affecting the relationship between smoking and lung cancer, we can treat the provided probabilities as the direct causal effects.} \\
\\ 
\texttt{1.  **Calculate the probability of lung cancer for a nonsmoker:**} \\
\texttt{    A nonsmoker can have either "no tar deposit" or "high tar deposit." We are given the marginal probability of smoking is 67\%, implying the probability of being a nonsmoker is \$1 {-} 0.67 = 0.33\$.} \\
\texttt{    *   \$P(\textbackslash{}text\{High Tar\} | \textbackslash{}text\{Nonsmoker\}) = 36\textbackslash{}\% = 0.36\$} \\
\texttt{    *   \$P(\textbackslash{}text\{No Tar\} | \textbackslash{}text\{Nonsmoker\}) = 1 {-} 0.36 = 0.64\$} \\
\texttt{    *   \$P(LC | \textbackslash{}text\{Nonsmoker, No Tar\}) = 12\textbackslash{}\% = 0.12\$} \\
\texttt{    *   \$P(LC | \textbackslash{}text\{Nonsmoker, High Tar\}) = 44\textbackslash{}\% = 0.44\$} \\
\\ 
\texttt{    Using the law of total probability:} \\
\texttt{    \$\$P(LC | \textbackslash{}text\{Nonsmoker\}) = (0.64 \textbackslash{}times 0.12) + (0.36 \textbackslash{}times 0.44)\$\$} \\
\texttt{    \$\$P(LC | \textbackslash{}text\{Nonsmoker\}) = 0.0768 + 0.1584 = 0.2352 \textbackslash{}text\{ (or \} 23.52\textbackslash{}\%)\$\$} \\
\\ 
\texttt{2.  **Calculate the probability of lung cancer for a smoker:**} \\
\texttt{    Similarly, for a smoker:} \\
\texttt{    *   \$P(\textbackslash{}text\{High Tar\} | \textbackslash{}text\{Smoker\}) = 84\textbackslash{}\% = 0.84\$} \\
\texttt{    *   \$P(\textbackslash{}text\{No Tar\} | \textbackslash{}text\{Smoker\}) = 1 {-} 0.84 = 0.16\$} \\
\texttt{    *   \$P(LC | \textbackslash{}text\{Smoker, No Tar\}) =} \\
            }
        }
    }
    & & \\ \\

    \theutterance \stepcounter{utterance}  
    & & & \multicolumn{2}{p{0.3\linewidth}}{
        \cellcolor[rgb]{0.9,0.9,0.9}{
            \makecell[{{p{\linewidth}}}]{
                \texttt{\tiny{[GM$|$GM]}}
                \texttt{Yes} \\
            }
        }
    }
    & & \\ \\

    \theutterance \stepcounter{utterance}  
    & & & \multicolumn{2}{p{0.3\linewidth}}{
        \cellcolor[rgb]{0.9,0.9,0.9}{
            \makecell[{{p{\linewidth}}}]{
                \texttt{\tiny{[GM$|$GM]}}
                \texttt{yes} \\
            }
        }
    }
    & & \\ \\

    \theutterance \stepcounter{utterance}  
    & & & \multicolumn{2}{p{0.3\linewidth}}{
        \cellcolor[rgb]{0.9,0.9,0.9}{
            \makecell[{{p{\linewidth}}}]{
                \texttt{\tiny{[GM$|$GM]}}
                \texttt{game\_result = WIN} \\
            }
        }
    }
    & & \\ \\

\end{supertabular}
}

\end{document}
